\chapter{Space Elevators and Expandable Structures}
\label{ch:24}

The space elevator is the oldest of the megascale concepts in this book and the one whose physics is simplest to state. A cable anchored on the equator and extending beyond geostationary altitude, with a counterweight at its far end, remains taut by the balance of gravitation and centrifugal force, and payloads can climb it without rockets. The idea originates with Tsiolkovsky~\cite{tsiolkovsky1895} in the form of a tower and with Artsutanov~\cite{artsutanov1960} and Pearson~\cite{pearson1975} in the form of a cable, and the analysis of Pearson in 1975 established the central quantitative fact, which is that the taper of the cable required to carry its own weight is governed by the specific strength of the material. The engineering study of Edwards and Westling~\cite{edwards2003} is the standard modern account of the cable. The author's notes pair the elevator with a second idea, the Hoberman-inspired expandable structure, which they use for nested towers that extend by unfolding, for orbital rings and tunnels, and as a lattice of the kind drawn over a globe in the physical models of the nine-circle arrangement. The notes do not give a mechanical model of what the expandable linkage contributes to the problem of bearing load, and the present chapter supplies one, with a conclusion that is unwelcome to the concept and is nonetheless the one the physics supports: expandable linkages solve the problem of packaging and deployment, and they do not solve the problem of strength. The chapter therefore separates the two roles and assesses each.

The elevator proper is a tension structure. At each radius $r$ along the cable the load it must carry is the net of the gravitational pull toward the Earth and the outward centrifugal force of the corotating frame, which is the derivative of the effective potential $\Psi(r)=-GM/r-\tfrac12\omega^2r^2$. The potential increases with radius up to the geostationary radius of $42{,}164$ km, where the net force vanishes, and the tension in a cable of uniform stress accumulates according to the difference of the potential between the surface and that point. Evaluating the difference gives $\Delta\Psi=4.84\times10^{7}$ J per kilogram. A cable designed so that the stress is the same at every height must have a cross-section that varies with height, and the ratio of its cross-section at geostationary altitude to that at the surface is the exponential of the ratio of $\Delta\Psi$ to the specific strength of the material. The formula leaves little room for argument. At the specific strengths, taking the full tensile strength as the working stress to give the material the most favorable treatment, of aramid, $2.5$ MJ per kilogram, the taper ratio is $e^{19.4}\approx2.6\times10^{8}$; of polybenzoxazole, $3.7$ MJ per kilogram, $e^{13.0}\approx4.5\times10^{5}$; of steel at a working stress of $0.7$ GPa, a number with more than two hundred digits; and of a hypothetical yarn with a strength of $50$ GPa and a density of $1300\ \mathrm{kg\,m^{-3}}$, $e^{1.26}\approx3.5$. A taper of ten to the fifth or more is not a design but a refutation, since the cross-section at the surface must carry the working load of a payload while the cross-section at the top must be a hundred thousand times larger. The elevator is therefore a concept that depends entirely on a material, available in macroscopic lengths of tens of thousands of kilometers with a strength of tens of gigapascals at densities near that of carbon, that has not been produced; and a safety factor of two squares the taper ratios of every entry above, so that even the hypothetical yarn rises to $12$ and the polybenzoxazole to $2\times10^{11}$. The reasons for the optimism of the concept's proponents are carbon nanotubes, whose individual strengths approach the required figure, and the reasons for the skepticism of its critics are that the strength of an assembly of nanotubes is limited by the weakest defect in a very long structure and by the imperfect transfer of load between the elements. The question of whether such a material can be made is a question of unresolved science and not of engineering, and the readiness class follows from it.

If the material were available, the remaining problems are of a different order but are not trivial. The cable would be exposed to debris and micrometeoroids along its length, in the region of low orbit where the debris is densest, and the survivability analysis of the preceding chapter applies. It would be subject to oscillation, since a tension structure that is tens of thousands of kilometers long has many modes, excited by the climbers, by the tides, and by the atmosphere at the base. It would require a moving platform at the base or a stabilizing counterweight at the top to manage the dynamics, and its traffic is limited by the power available to the climbers. A climber that moves at $200\ \mathrm{km/h}$ takes $7.5$ days to reach the geostationary station, and a payload of one tonne at the surface requires a mechanical power of about $0.54$ MW to maintain that speed, declining as the effective gravity falls with height. The energy to be supplied per kilogram is the difference of the potential, $48.4$ MJ or $13.5$ kWh, which is larger than the energy of orbital velocity in low orbit because the destination is higher; the elevator does not reduce the energy required to reach geostationary altitude but allows it to be supplied from the ground, by electricity, as in the geothermal arrangement of Chapter 21, rather than by the combustion of propellant.

The second half of the pairing concerns the structures of the Hoberman type. A scissor linkage consists of pairs of bars pinned at their midpoints, so that the bars can be rotated together from a compact folded state to an extended state, and the Hoberman sphere, devised by Chuck Hoberman in the 1990s, is an arrangement of such linkages in which a sphere expands uniformly to several times its folded diameter. The geometric property is the ratio between the stowed and deployed dimensions. For a stack of scissor cells whose bars make an angle $\theta$ with the horizontal, the height of each cell is proportional to $\sin\theta$, and the expansion ratio between a stowed angle of five degrees and the fully deployed angle of ninety degrees is $1/\sin5^\circ=11.5$; for a stowed angle of three degrees it is $19$. Nested or telescoping arrangements, in which successive stages are stowed inside one another and extend sequentially, multiply the achievable extension by the number of stages. These are real and useful properties, which suit the structure to the transport of a large volume in a small package, as is required for anything built in orbit from components lifted by rockets, and they suit it for adjusting its dimensions in response to loads. They are the properties by virtue of which expandable structures have found their uses on spacecraft in the form of antennas, solar arrays, and booms.

The load-bearing function is another matter, and here the physics is conclusive. A locked scissor lattice is a truss with additional joints, and its strength in compression is limited by the buckling of its members and by the stiffness of its joints, which are less than those of a welded or bonded truss of the same mass. A tower is a column, and a column is limited by the self-buckling length computed by Greenhill in 1881~\cite{greenhill1881}: a uniform column of bending stiffness $EI$ and weight per unit length $q$ buckles under its own weight when its height exceeds $(7.84\,EI/q)^{1/3}$. For a thin-walled tube of radius $r$ this gives a limiting height of $1.58\,(E/\rho g)^{1/3}r^{2/3}$, which is proportional to only the two-thirds power of the radius, so that large height demands radius out of all proportion. For steel, a tube of radius $5$ m can stand $630$ m, of radius $50$ m $2.9$ km, and of radius $500$ m $13.6$ km. To reach $20$ km a steel tube would need a radius of $890$ m, and to reach $100$ km, ten kilometers. With a material of the modulus of an ideal carbon nanostructure, one terapascal at a density of $1300\ \mathrm{kg\,m^{-3}}$, the radii are $160$ m for $20$ km and $1.8$ km for $100$ km. A self-supporting compression tower to the edge of space is therefore excluded for any material of known modulus, whatever the geometry of its joints; an actual tall structure is tapered, guyed, or braced against a wide base, which changes the numbers by a modest factor but not the conclusion, and the tallest structures built are under a kilometer. The proposals in the notes in which nested expandable towers extend to great heights, and carry within them linear ring tunnels for launching, must therefore be interpreted either as structures of tens of kilometers whose feasibility would depend on pneumatic or tensile support and has not been analyzed, or as structures assembled in orbit, where gravity does not load them in the same way. In orbit, the expandable lattice comes into its own: a ring or habitat that must be built from components lifted separately can be folded into a launch vehicle and deployed on station, and the strength it requires is that needed to hold its own shape against centrifugal and tidal loads, which is modest.

The ordering that the author has proposed for the use of expandable structures, which is to build first a ring with compressive strength of its own so that the scissor-like motion is confined to the macroscopic act of deployment, is the correct one, and it is the ordering that the argument above recommends. In it the load-bearing member is a rigid structure whose strength does not depend on its joints, and the expandable linkage serves as the means of carrying the structure into place and opening it, after which it is locked and plays no further part in bearing load. The question that remains is what loads the ring itself must bear, and the answer depends sharply on how the ring moves. A thin ring of radius $r$ that turns about the Earth's center at an angular rate $\Omega$ has a hoop stress equal to the density times the difference between the square of its own speed and the square of the local circular orbital speed, $\sigma=\rho(v^2-v_{\mathrm{orb}}^2)$: positive, a tension, if the ring moves faster than a circular orbit, zero if it moves at orbital speed, and negative, a compression, if it moves slower. A ring in free fall at orbital speed therefore bears no hoop load at all, and this is the condition in which an expandable ring is most easily built, deployed in orbit from components lifted separately and opened to its full radius. A ring that is to be stationary with respect to the ground, however, moves slowly and is crushed: at low altitude the compressive stress is the density times roughly $5.9\times10^{7}\ \mathrm{J\,kg^{-1}}$ at $400$ km altitude and $6.2\times10^{7}$ at the surface, which for steel is nearly $490$ GPa and which for any material requires a compressive specific strength greater than that of the strongest known fibers in tension. The figure falls to $3.9\times10^{7}$ at a radius of $10{,}000$ km, to $1.8\times10^{7}$ at $20{,}000$ km, and to $8.5\times10^{6}$ at $30{,}000$ km, and it vanishes at the geostationary radius, where a corotating ring is in free fall as well, and beyond which it becomes a tension. The geostationary ring is thus the one place at which a ring that stays fixed above the ground bears no hoop load, and it is where the elevator's cable is balanced, so that the two structures meet naturally there. The conclusion for the ring-first architecture is that a rigid ring should be built in free fall or at geostationary radius, where the compressive requirement is small, and that a rigid ring that is fixed close to the ground is excluded by the same arithmetic that excludes the tower. It also follows that the ring is a hub and not a road: a ring in low orbit moves past the ground at nearly eight kilometers per second, and the connection to the ground must be made by the skyhooks and accelerators of the preceding chapters, or by cables that descend from a ring at geostationary altitude, which return the argument to the tensile material on which the elevator depends. The scale of a ring is also considerable: its circumference at $400$ km altitude is $42{,}500$ km, and even at a mass of one kilogram per meter it contains $4.3\times10^{4}$ tonnes, at one hundred kilograms per meter $4.3$ million tonnes. A free-fall ring is, moreover, subject to Maxwell's instability, so that it must be actively stabilized, and the cost of that stabilization is a continuing obligation. The ring-first architecture is thus a coherent program of orbital construction, and it is placed in the second readiness class as such, while its connection to the ground by anything but skyhooks and accelerators remains in the third.

The assessment under the standard of Chapter 16 is that the elevator is in the third readiness class until a material exists, that the expandable linkage is in the first class for packaging and the second for orbital assembly of large structures, and that the combination of the two as a tower to space is excluded by the argument above. The book records this as an impossibility and not as a speculative possibility, in keeping with its commitment to state the failures of proposals as carefully as their promise. The hub structures that serve in the notes as landing places for the expandable towers, the intervolsorial pediments of Chapter 20, can be treated separately as large offshore platforms, which is a feasible engineering category, and the question of what arrives on them is the question of the preceding chapters. What remains of the concept after the exclusion is a program of materials science, directed at the production of high-strength fiber in macroscopic lengths, which is worth pursuing for its own sake, and of orbital construction, which is the subject of the following chapters.

\section*{Formalization}

Two results are given, the first for tension structures and the second for compression columns.

Let a cable of density $\rho$ have working stress $\sigma$ at every cross-section, and let $A(r)$ be the cross-sectional area at geocentric radius $r\in[R,r_g]$, with $\Psi(r)=-GM/r-\tfrac12\omega^2r^2$ the effective potential of the corotating frame and $r_g=(GM/\omega^2)^{1/3}$ the geostationary radius.

\begin{proposition}
Equilibrium of a cable element requires $\sigma\,A'(r)=\rho A(r)\,\Psi'(r)$, so $A(r)=A(R)\exp\!\big[(\rho/\sigma)(\Psi(r)-\Psi(R))\big]$. Since $\Psi'(r)=GM/r^2-\omega^2r>0$ for $r<r_g$, the area increases monotonically to $r_g$, where $\Psi'=0$. The taper ratio $A(r_g)/A(R)=\exp(\Delta\Psi\,\rho/\sigma)$ with $\Delta\Psi=\Psi(r_g)-\Psi(R)$. A climber of mass $m$ moving at speed $u$ at radius $r$ requires mechanical power $mu\,\Psi'(r)$, and the energy to ascend from $R$ to $r_g$ is $m\Delta\Psi$.
\end{proposition}

\begin{proof}
For the cable element between $r$ and $r+dr$, the tension at the outer end exceeds that at the inner end by the net of the gravitational and centrifugal forces on the element, which is $-\rho A\Psi'\,dr$ directed outward-minus-inward, giving $d(\sigma A)=\rho A\Psi'\,dr$ with $\sigma$ constant. Integrate. The climber's force against the effective gravity $\Psi'$ times speed gives the power, and integrating over $r$ gives the energy.
\end{proof}

With $GM=3.986004\times10^{14}\ \mathrm{m^3\,s^{-2}}$, $\omega=7.2921\times10^{-5}\ \mathrm{s^{-1}}$ and $R=6378$ km, one obtains $r_g=42{,}164$ km, $\Delta\Psi=4.842\times10^7\ \mathrm{J\,kg^{-1}}$ ($13.45$ kWh kg$^{-1}$), and the taper exponent $\Delta\Psi\rho/\sigma$ equal to $543$, $19.4$, $13.0$, $1.26$, and $0.63$ for working-stress-to-density ratios of $0.089$ (steel at $0.7$ GPa), $2.5$ (aramid), $3.7$ (polybenzoxazole), $38.5$ ($50$ GPa at $1300$ kg m$^{-3}$), and $76.9$ MJ kg$^{-1}$ ($100$ GPa). The surface effective gravity is $9.78\ \mathrm{m\,s^{-2}}$, giving climber power $0.54$ MW per tonne at $55.6\ \mathrm{m\,s^{-1}}$, and the ascent time at that speed is $7.5$ days.

For the compression column, let a uniform thin-walled tube of radius $r_0$ and wall thickness $t\ll r_0$ stand with its base clamped and its top free. The weight per unit length is $q=2\pi r_0t\rho g$ and the second moment of area is $I=\pi r_0^3t$.

\begin{proposition}
The column buckles under its own weight at the critical height $H_{\mathrm{cr}}=(\lambda_1EI/q)^{1/3}$, with $\lambda_1\approx7.84$ (Greenhill). For the tube, $H_{\mathrm{cr}}=(\lambda_1/2)^{1/3}\,(E/\rho g)^{1/3}\,r_0^{2/3}\approx1.58\,(E/\rho g)^{1/3}r_0^{2/3}$, independent of the wall thickness. To reach a height $H$ requires $r_0=\big[H/(1.58(E/\rho g)^{1/3})\big]^{3/2}$.
\end{proposition}

The substitution is immediate: $EI/q=Er_0^2/(2\rho g)$. The value of $\lambda_1$ is the first root of the Bessel-function eigenvalue condition for the self-buckling of a clamped-free column and is quoted without derivation. Numerically, for steel ($E=200$ GPa, $\rho=7850$) $H_{\mathrm{cr}}=633$ m, $2.94$ km, $13.6$ km at $r_0=5$, $50$, $500$ m, and $r_0=887$ m and $9.9$ km for $H=20$ km and $100$ km; for a $1$ TPa, $1300\ \mathrm{kg\,m^{-3}}$ material $H_{\mathrm{cr}}=1.97$, $9.2$, $42.5$ km at the same radii, and $r_0=161$ m and $1.8$ km for $H=20$ and $100$ km. The $2/3$ power of the dependence on radius means that increasing height tenfold requires increasing radius by a factor $10^{3/2}=31.6$, which is the quantitative content of the exclusion of compression towers to the edge of space.

For the expandable lattice, a scissor cell whose bars of total length $2b$ cross at their midpoints at an angle $\theta$ to the horizontal has height $2b\sin\theta$, so a stack of $n$ cells has height $H(\theta)=2nb\sin\theta$ and expansion ratio $H(\pi/2)/H(\theta_s)=1/\sin\theta_s$, equal to $19.1$, $11.5$, $5.8$, $2.9$ for $\theta_s=3^\circ,5^\circ,10^\circ,20^\circ$. The ratio depends on the stowed angle alone and is independent of the strength of the bars, which is the sense in which the expansion ratio and the load capacity are separate properties of the structure.

For the ring-first architecture, consider a thin ring of density $\rho$ and radius $r$ moving tangentially at speed $v=\Omega r$ about the Earth's center, in the plane of the equator.

\begin{proposition}
The hoop stress in the ring, positive in tension, is $\sigma=\rho\,(v^2-GM/r)$. It vanishes for $v=v_{\mathrm{orb}}=\sqrt{GM/r}$ (free fall), is compressive for $v<v_{\mathrm{orb}}$ with magnitude $\rho(v_{\mathrm{orb}}^2-v^2)$, and is tensile for $v>v_{\mathrm{orb}}$. For a ring corotating with the Earth, $v=\omega r$, the stress is $\sigma=\rho\,(\omega^2r^2-GM/r)$, which is compressive for $r<r_g$, zero at $r=r_g=(GM/\omega^2)^{1/3}$, and tensile for $r>r_g$. A ring of compressive specific strength $s_c$ can be corotating at radius $r$ only if $s_c\ge GM/r-\omega^2r^2$.
\end{proposition}

\begin{proof}
A ring element of length $r\,d\varphi$ and cross-section $A$ has mass $\rho A\,r\,d\varphi$. Its centripetal requirement is $\rho Arv^2/r\,d\varphi$ directed inward, supplied by gravity $\rho A\,r\,d\varphi\,GM/r^2$ and by the net inward component of the hoop force $2N\sin(d\varphi/2)\approx N\,d\varphi$. Hence $N=\rho A(v^2-GM/r)$, and $\sigma=N/A$.
\end{proof}

Numerically, the required compressive specific strength $GM/r-\omega^2r^2$ is $62.3$ MJ kg$^{-1}$ at $r=6378$ km, $58.6$ at $6771$ km, $39.3$ at $10{,}000$ km, $17.8$ at $20{,}000$ km, $8.5$ at $30{,}000$ km, and $0$ at $42{,}164$ km. For steel at $r=6378$ km the stress is $7850\times6.23\times10^7=489$ GPa. For comparison, the specific strengths of the strongest fibers in tension are of the order of $4$ MJ kg$^{-1}$ (polybenzoxazole) to a few tens of MJ kg$^{-1}$ (hypothetical nanotube yarns), and compressive specific strengths are lower. The circumference of a ring at $6771$ km radius is $42{,}543$ km, so its mass at linear density $\lambda$ is $4.25\times10^{4}\lambda$ kg, which is $4.3\times10^{4}$ tonnes for $\lambda=1$ kg m$^{-1}$ and $4.3\times10^{6}$ tonnes for $\lambda=100$ kg m$^{-1}$. A planar circular scissor mechanism that deploys the ring from a stowed state has radial expansion ratio $1/\sin\theta_s$ for the same reason as the stack of the preceding paragraph, $11.5$ at $\theta_s=5^\circ$, and for this reason a ring that is to be deployed to a radius of thousands of kilometers must either be built in segments and joined or be a ring whose stowed form is itself a large structure; in neither case does the linkage carry the loads of the deployed ring.
